%قبل از شروع مسئله ترجیحاً چیزی ننویسید. اما اگر خواستید بنویسید، در قالب لیستی (فقرات و شمارش) بنویسید.
%برای اضافه کردن نمره، بعدی از شروع‌مسئله، [x] را اضافه کنید و به جای x  نمره مورد نظر را قرار دهید. برای زیرمسئله‌ها هم به همین صورت است.
%\شروع‌مسئله[10]
\شروع‌مسئله [20]
\متن‌سیاه{زوج مرتب های مرتب}

فرض کنید $1,000,000$ زوج مرتب مثل $(a,b)$ داریم که باید آن ها را طبق قانون زیر مرتب کنیم:
$$(a_1,b_1) < (a_2,b_2) \iff (a_1 < b_1)\lor\big( (a_1=b_1) \land (b_1<b_2)\big)$$


\متن‌سیاه{ورودی:}
شامل $1,000,000$ خط است که در هر خط زوج مرتب $(a,b)$ به ترتیب به کاربر داده می شود.
\متن‌لاتین{
	\begin{center}
		$1 \le a < 100\,000$
		\linebreak[1]
		$1 \le b \le 1\,000\,000\,000\,000\,000$
	\end{center}
}
\متن‌سیاه{خروجی:}
باید زوج مرتب‌ها پس از مرتب کردن باشد (در هر خط یک زوج مرتب).

\متن‌سیاه{نکات:}
در این مسئله شما تنها مجاز به استفاده از زبان پایتون و جاوا هستید.

ورودی دقیقا $1,000,000$ زوج مرتب دارد اما به علت محدودیت صورت سؤال، مثال زیر برای 10 زوج مرتب است.

\متن‌سیاه{ورودی و خروجی نمونه:}
\begin{latin}
	\begin{tabularx}{\textwidth}{|X | X|}
		\hline
		\textbf{input} &
		\textbf{output}\\
		\hline
		13241 323294848907982&
		13241 152840012736194\\
		
		97255 553012130624557 &
		13241 241855619495066\\
		
	    13241 876302146194655 &
		13241 323294848907982\\
		
	    33481 299942243878034&
		13241 876302146194655\\
		
		97255 858692723362742 &
		33135 358707532837541\\
		
		33135 358707532837541 &
		33481 299942243878034\\
		
		13241 241855619495066 &
		49483 395372516597937\\
		
			63663 186195109346510 &
	63663 186195109346510\\
		
				49483 395372516597937 &
		97255 553012130624557\\
		
				13241 152840012736194 &
		97255 858692723362742\\
		\hline
		
	\end{tabularx}
\end{latin}


%اگر مسئله چند قسمت داشت به این صورت عمل کنید:
%\شروع{زیرمسائل}
%\زیرمسئله[0] % مسئلۀ ۱آ
%\زیرمسئله % مسئلۀ ۱ب
%\زیرمسئله % مسئلۀ ۱پ
%\زیرمسئله % مسئلۀ ۱ت
%\پایان{زیرمسائل}

\solution{
مقدار $a$ محدود است، پس می‌توانیم بر اساس زوج مرتب‌ها را در براساس $a$هایشان در سطل‌های مختلف بریزیم و اعداد داخل هر سطل را با یک الگوریتم مرتب سازی دلخواه مرتب کنیم (مثل مرتب‌سازی شمارشی یا ادغامی). کد آقای وحید زهتاب با نام \متن‌تایپ{ordered-pair.py} به عنوان نمونه در پوشه مربوطه قرار داده شده.
	
}




%%%%%%%%%%%%%%%%%%%%%%%%%%%%%%%%%%%%%%%%%%%%%%%%%%%%%
% برخی دستورات مفید برای تمرین‌ها. %
%%%%%%%%%%%%%%%%%%%%%%%%%%%%%%%%%%%%%%%%%%%%%%%%%%%%%

% فرمول‌های نمونه:
%$f(x) = \Theta(x)$
%$T(x, y) \leq \log(x) + 2^y + \binom{2n}{n}$
% {\tt code\_function}


% لیست موارد بدون شماره‌گذاری:
%\شروع{فقرات}
%    \فقره مورد اول 
%        \شروع{فقرات}
%            \فقره مورد اول داخلی 
%                \شروع{فقرات}
%                    \فقره اولین مورد خیلی داخلی 
%                    \فقره دومین مورد خیلی داخلی 
%                \پایان{فقرات}
%            \فقره مورد دوم داخلی 
%        \پایان{فقرات}
%    \فقره مورد دوم 
%\پایان{فقرات}

% لیست موارد شماره‌گذاری شده:
%\شروع{شمارش}
%    \فقره یک مورد 
%    \فقره یک مورد دیگر 
%    \فقره باز هم یک مورد دیگر
%\پایان{شمارش}

% برای نوشتن یک فرمول چند خطی که یک شماره بخوره. برای حالت‌های دیگر به سایت زیر مراجعه کنید:
% https://www.overleaf.com/learn/latex/Aligning_equations_with_amsmath
%\begin{align} 
%    \begin{split}
%        (x+y)^3 &= (x+y)^2(x+y) \\
%                &= (x^2+2xy+y^2)(x+y) \\
%                &= (x^3+2x^2y+xy^2) + (x^2y+2xy^2+y^3) \\
%                &= x^3+3x^2y+3xy^2+y^3
%    \end{split}                                 
%\end{align}

% نوشتن ماتریس و آرایه دوبعدی
%\begin{align}
%    A = 
%    \begin{bmatrix}
%        A_{11} & A_{21} \\
%        A_{21} & A_{22}
%    \end{bmatrix}
%\end{align}

% اضافه کردن تصویر با شرح. برای اضافه کردن تصویر بدون شرح، فقط دستور درج‌تصویر را استفاده کنید.
%\شروع{شکل}[h!]
%\تنظیم‌ازوسط
%  \درج‌تصویر[پهنا=0.5\پهنای‌متن]{img.jpg}
%\شرح{شرح تصویر}
%\برچسب{تصویر:۱}
%\پایان{شکل}
%\پایان{زیرمسائل}

%کد زیر را ببینید.
%%\lstinputlisting{source_filename.py}
%\begin{latin}
%\begin{lstlisting}
%# Import the modules
%import sys
%import random
%
%ans = True
%
%while ans:
%question = raw_input("Ask the magic 8 ball a question: (press enter to quit) ")
%
%answers = random.randint(1,8)
%
%if question == "":
%sys.exit()
%
%elif answers == 1:
%print "It is certain"
%
%elif answers == 2:
%print "Outlook good"
%
%elif answers == 3:
%print "You may rely on it"
%
%elif answers == 4:
%print "Ask again later"
%
%elif answers == 5:
%print "Concentrate and ask again"
%
%elif answers == 6:
%print "Reply hazy, try again"
%
%elif answers == 7:
%print "My reply is no"
%
%elif answers == 8:
%print "My sources say no"
%\end{lstlisting}
%\end{latin}
